Die Tabelle \emph{Thermische Daten von Gasen} (FoTaBe S.~191) gibt für Kohlendioxid $c_p = \cpCO$; die molare Masse von CO$_2$ ist $M = \MCO$ (FoTaBe S.~228). Die molaren Wärmen sind
\begin{align}
 C_p &= \CpCF = \cpC\cdot\MC = \CpC \\
 C_v &= C_p - R = \CvC
\end{align}
Aus $C_v = \frac{f}{2}R$ folgt
\begin{align}
 f &= \frac{2\,C_v}{R} = \fCOF \\
   &= \frac{2\cdot\CvC}{\ncR} = \fCO \approx \result{\fCOP{0}{2}}
\end{align}
Kohlendioxid hat also etwa sieben angeregte Freiheitsgrade: drei der Translation, zwei der Rotation (das Molekül ist linear) und zusätzlich teilweise angeregte Schwingungen.
